\section{Factorization Current Conventions on the Brane Line}
\label{sec:app-factorization-current-conventions}

This appendix fixes the current conventions used in
Theorem~\ref{thm:lane-factorization-current}.  Its stable core is the
degree-zero Hamiltonian current source, the continuous current dual,
the CE coordinate which maps to \(B_f(a)\), and the density sign under
orientation reversal.  It does not construct the unreduced BV
cotangent factorization centre; that remains the obligation recorded
in Remark~\ref{rmk:unreduced-lift-status}.

\begin{defn}[Hamiltonian current source]
\label{def:app-factorization-current-source}
  For an open interval \(I\subset\R\), the reduced Hamiltonian current
  Lie algebra is
  \[
    \mathfrak h_I=\Omega^\bullet_c(I)\widehat\otimes\mathfrak h .
  \]
  In the polynomial degree-zero formulas one may replace
  \(\mathfrak h\) by its dense scalar-reduced polynomial subalgebra
  \(\bar A=\C[z_1,z_2]/\C\cdot1\).  If
  \(\alpha,\beta\in\Omega^\bullet_c(I)\) and \(f,g\in\mathfrak h\),
  the bracket is
  \[
    [\alpha\otimes f,\beta\otimes g]
      =\alpha\wedge\beta\otimes\{f,g\}.
  \]
  For degree-zero test functions \(a,b\in\Omega^0_c(I)\) this is
  \[
    [a\otimes f,b\otimes g]=ab\otimes\{f,g\}.
  \]
  The measurement notation writes the oriented density associated to
  \(a\) as \(a(t)\,dt\); this density notation does not change the
  de Rham degree of the source test function \(a\).
\end{defn}

\begin{defn}[Current dual and shifted cotangent source]
\label{def:app-factorization-current-dual}
  Let
  \[
    \mc P\simeq\mathfrak h^\vee_{\mathrm{cont}}
  \]
  be the reduced Matlis principal-part module of
  Section~\ref{sec:app-matlis-principal-parts}.  The current dual used
  on \(I\) is
  \[
    \mathcal P_I=\mathcal D'_c(I)\widehat\otimes\mc P,
  \]
  with pairing
  \[
    \langle B\otimes\rho,\;a\otimes f\rangle
      =B(a)\operatorname{Res}(\rho f).
  \]
  If \(B=b(t)\,dt\) is a regular compactly supported density, then
  \(B(a)=\int_I a(t)b(t)\,dt\).

  The interval shifted-cotangent Lie algebra is
  \[
    \mathfrak g_I=\mathfrak h_I\ltimes\mathcal P_I[1].
  \]
  In degree zero the shifted cotangent generator
  \(B\otimes\rho[1]\) is odd.  The CE coordinate dual to such a
  shifted cotangent generator has degree zero.  We denote the
  coordinate dual to the cotangent functional labelled by
  \(a(t)\,dt\otimes f\) by
  \[
    u_{a(t)dt\otimes f},\qquad |u_{a(t)dt\otimes f}|=0.
  \]
  The Hamiltonian CE coordinate \(c_{a\otimes f}\), dual to the
  unshifted Hamiltonian generator \(a\otimes f\), has degree one.
\end{defn}

\begin{prop}[Source of the smeared boundary Hamiltonian]
\label{prop:app-factorization-source-of-B}
  The length-one shifted-cotangent CE coordinate, not a shifted dual of
  the unshifted Hamiltonian current, is the source of the smeared
  Hamiltonian operation:
  \[
    u_{a(t)dt\otimes f}\longmapsto
    B_f(a)
      =\int_I a(t)\,
        \overline{\operatorname{Tr}
        f(\phi_1(t),\phi_2(t))}\,dt .
  \]
  With this convention the degree-zero current bracket is
  \[
    \{B_f(a),B_g(b)\}=B_{\{f,g\}}(ab).
  \]
\end{prop}

\begin{proof}
  The CE algebra of
  \(\mathfrak g_I=\mathfrak h_I\ltimes\mathcal P_I[1]\) is the
  completed symmetric algebra on the continuous dual shifted by
  \([-1]\).  A generator of \(\mathcal P_I[1]\) has degree one in the
  degree-zero current sector; its CE coordinate therefore has degree
  zero and is written \(u\).  Under the CE/PV dictionary this
  coordinate maps to multiplication by the corresponding Hamiltonian
  boundary observable.  Applying boundary evaluation
  \(f\mapsto\overline{\operatorname{Tr}f(\phi_1,\phi_2)}\) and
  smearing against the compactly supported density \(a(t)\,dt\) gives
  the displayed \(B_f(a)\).  The bracket identity is then the
  equal-time Poisson bracket smeared against \(a(t)b(s)\delta(t-s)\),
  which gives the pointwise product \(ab\).
\end{proof}

\begin{prop}[Principal-part current bracket]
\label{prop:app-factorization-principal-part-bracket}
  In the reduced \(P_0\) current model after de Rham-current contraction,
  \[
    A_\partial^{\mathrm{pp}}(I)
      =\widehat{\mathbf S}(\mathfrak h_I)
        \widehat\otimes
        \widehat{\mathbf S}(\mathcal P_I[1]),
  \]
  the degree-zero brackets are
  \[
  \begin{aligned}
    \{B_f(a),B_g(b)\}&=B_{\{f,g\}}(ab),\\
    \{B_f(a),\Theta_\rho(B)\}&=\Theta_{f\cdot\rho}(aB),\\
    \{\Theta_\rho(B),\Theta_\sigma(C)\}&=0.
  \end{aligned}
  \]
  Here \(aB\) is multiplication of a compactly supported current by a
  smooth test function.  On monomial labels and a regular current
  \(B=b(t)\,dt\),
  \[
    \{B_{z_1^p z_2^q}(a),\Theta_{\rho_{i,j}}(B)\}
      =-(pj-qi+p-q)\Theta_{\rho_{i-p+1,j-q+1}}(aB),
  \]
  with negative-index targets equal to zero.
\end{prop}

\begin{proof}
  The first bracket is the Hamiltonian current bracket from
  Proposition~\ref{prop:app-factorization-source-of-B}.  The second
  bracket is the negative transpose of the first under the product of
  the de Rham-current pairing and the residue pairing:
  \[
    \langle B\otimes\rho,\;a\otimes f\rangle
      =B(a)\operatorname{Res}(\rho f).
  \]
  The residue factor gives the coadjoint action
  \(f\cdot\rho\), computed in
  Proposition~\ref{prop:app-matlis-coadjoint-action}; the line factor
  gives \(aB\).  The \(\Theta\)-sector is declared abelian in the
  reduced semidirect product, giving the third bracket.
\end{proof}

\begin{rmk}[Orientation reversal]
\label{rmk:app-factorization-orientation-reversal}
  The orientation convention is the ordinary one for one-dimensional
  densities.  Under an orientation-reversing coordinate change,
  \(dt\) contributes its usual sign.  No additional sign is inserted
  from the CE coordinate \(u_{a(t)dt\otimes f}\), from the
  principal-part label \(\rho\), or from the shifted cotangent
  notation.  Thus the displayed operations are orientation-covariant
  after transforming the densities and currents in the usual way; if
  one reverses only the chosen orientation while holding the written
  density fixed, the sign is exactly the sign of \(dt\).
\end{rmk}

\begin{rmk}[Residual convention boundary]
\label{rmk:app-factorization-residual}
  These conventions prove the reduced \(P_0\) current bookkeeping.
  They do not supply the missing unreduced BV current complex, the
  primitive cumulant projection before reduction, or the homotopies
  making \(\Theta_\rho\) central against arbitrary unreduced open
  observables.
\end{rmk}
